%%%%%%%%%%%%%%%%%%%%%%%%%%%%%%%%%%%%%%%%%%%%%%%%%%%%%%%%%%%%%%%%%%%%%%%%%%%%%%%
\section{Results on simulations}
\label{sec:simres}

The simulation set has \nmSimCases\ cases on \nmSimStars\ stars with a median of
\nmSimMedianN\ nightly epochs: \nmNPlanetOnly\ planets alone, \nmNPlanetAct\ planets with
activity, \nmNActOnly\ activity alone and \nmNNoiseOnly\ noise alone.

\subsection{False alarms}
\label{sec:fap}

For the false-alarm calibration we drew \nmFapPerStar\ noise-only series per star
(\nmFapSeries\ in all) of two kinds: the star's own velocities permuted within each
era, and Gaussian noise with the star's error bars and fitted jitters. Each went
through the same clip, trend and detection step as the simulations
(Fig.~\ref{fig:fapcal}). Neither common treatment of the jitter gives nominal rates.
With the jitter held at its planet-free value, the statistic of our first version, the
Baluev FAP of the highest peak falls below 1\% in only \nmFapPermFixOne\% of the
real-noise series (95\% interval \nmFapPermFixOneLo--\nmFapPermFixOneHi\%) and below
10\% in \nmFapPermFixTen\%: conservative by a factor of
\nmFapFixFactorLo--\nmFapFixFactorHi\ between 1 and 10\%, because a peak's own power
inflates the jitter that weights it. With the jitter rescaled under the signal model, as
in the ladder, the approximation errs the other way: \nmFapPermLROne\% of real-noise
series fall below 1\% (\nmFapPermLROneLo--\nmFapPermLROneHi\%), \nmFapPermLRTen\% below
10\% and \nmFapPermLRTenth\% below 0.1\%, a factor of \nmFapLiberalPerm\ at 1\% and
\nmFapLiberalPermTenth\ at 0.1\%. For Gaussian noise the factor at 1\% is
\nmFapLiberalGauss, so part of the excess comes from fitting the jitter scale, which the
known-variance approximation does not count, and part from the heavy tails of real
velocities. On these stars a Baluev FAP of \nmFapEmpOne\% (95\% interval
\nmFapEmpOneLo--\nmFapEmpOneHi\%) corresponds to a real false-alarm rate of 1\%, and
\nmFapEmpTenth\% (\nmFapEmpTenthLo--\nmFapEmpTenthHi\%) to 0.1\%. We keep the frozen 1\%
threshold in the rule and report its real rate. Permutation makes the noise white
between nights, so none of these rates includes correlated noise on timescales longer
than a night.

\begin{figure}
\centering
\includegraphics[width=\columnwidth]{fig_fapcal.pdf}
\caption{False-alarm calibration: the fraction of noise-only series whose highest
peak has a Baluev FAP below a given threshold, for the likelihood-ratio statistic of
the ladder (blue) and for $\dchi$ with the jitter fixed at its planet-free value
(red), on the stars' own permuted velocities (solid, with 95\% intervals) and on
Gaussian noise (dotted). The dashed line is the nominal rate.}
\label{fig:fapcal}
\end{figure}

\subsection{Completeness}

Planets alone are recovered at the right period with FAP below 1\% in half the cases
at an expected $K/\sigma_K$ of \nmCompFifty\ and in 90\% at \nmCompNinety\
(Fig.~\ref{fig:completeness}), with \nmCompPLo--\nmCompPHi\% recovery across five period
ranges between 1.2 and 300\,d at expected $K/\sigma_K > 10$. Of the \nmPlDetN\ series
with a planet alone and a significant highest peak, \nmPlWrongPct\% (\nmPlWrongN) peak at
the wrong period: \nmPlWrongAliasN\ at an alias or harmonic of the planet and the rest
on a noise peak that outgrew a weak planet.

\begin{figure}
\centering
\includegraphics[width=\columnwidth]{fig_completeness.pdf}
\caption{Recovery of injected planets (alone) at the right period with FAP below
1\%, against the expected $K/\sigma_K$.}
\label{fig:completeness}
\end{figure}

\subsection{What the impostors are}

Of the \nmNDet\ cases with a significant highest peak, \nmNDetPlanet\
(\nmDetPlanetPct\%) are the injected planet and \nmNDetImp\ are impostors:
\nmImpRotHarm\ at a rotation harmonic or its alias, \nmImpRotLine\ inside a line of the
activity spectrum but more than one resolution element from a harmonic (\nmImpEnvelopeN\
of them in the low-frequency envelope), \nmImpActOther\ elsewhere in series with
activity, \nmImpPlAlias\ at an alias or harmonic of the planet, and \nmImpNoise\ from
noise. These proportions follow from the class mix and the amplitude priors, not from
the archive. Because both planet and activity amplitudes scale with the star's scatter,
the planet fraction hardly depends on $N$ (\nmDetPlanetPctNLo\% in the lowest quartile,
\nmDetPlanetPctNHi\% in the highest), but it would change with the mix. Less dependent on
the mix is where activity impostors fall. An evolving spot pattern has lines of width
$\sim 1/(\pi\lambda)$, and of the \nmImpActTotal\ activity impostors \nmImpRotLinePctAct\%
lie inside those lines but more than one resolution element from any exact harmonic,
about as many as lie on a harmonic (\nmImpRotLineNoEnv\ of them excluding the
low-frequency envelope). The balance depends on the spot lifetime, and so on its prior:
for $\lambda$ shorter than three rotations, \nmLineShort\ impostors lie inside a line at
$k = 1$--4 and \nmHarmShort\ on a harmonic; for $\lambda$ longer than ten rotations,
\nmLineLong\ and \nmHarmLong. A harmonic-distance test does not look inside the lines.

\subsection{Calibration of the rungs on planets}
\label{sec:calib}

For the \nmNPlanetOnlyDet\ planets alone detected at the right period, the block phase
probability has a median of \nmBPhMed\ and falls below 0.05 in \nmBPhFive\% of cases,
close to nominal. With $n_{\rm b} - 2$ degrees of freedom for the phase, our first choice
on the argument that the fitted frequency absorbs a linear phase drift, \nmBPhFiveNbTwo\%
of the same planets fall below 0.05, so we adopted $n_{\rm b} - 1$. The block amplitude
probability is somewhat liberal: its median is \nmBAmpMed\ and \nmBAmpFive\%
(\nmBAmpFiveLo--\nmBAmpFiveHi\%) fall below 0.05, probably because an eccentric orbit
fitted with a sinusoid, and heavy-tailed noise, add scatter the block uncertainties do
not include. The half-phase difference exceeds
$1.96\sigma$ in \nmHPhOver\% of planets and the half-amplitude difference in
\nmHAmpOver\%. Taken alone, the rungs of the a-priori rule reject these planets at the
following rates: period \nmRungRejPeriod\%, phase \nmRungRejPhase\%, indicators
\nmRungRejInd\% and rotation \nmRungRejRot\%. The indicator rung, nominally a
$10^{-3}$ test, rejects \nmIndOverPl\% because it takes the maximum over indicators and
frequencies.

\subsection{Which tests carry the information}

Table~\ref{tab:auc} gives the power of each feature alone to separate planets from
impostors among significant peaks, as the area under the ROC curve (AUC).
No single test is decisive. The strongest, the indicator power (\nmAucInd), the block
phase rms (\nmAucBlockRms) and the growth of $\dchi$ (\nmAucGrowth), are only slightly
better than the detection statistics themselves (FAP \nmAucFap, $\dchi$ \nmAucDchi).
The distance to a rotation harmonic is informative where it can be computed: for the
\nmHarmDefinedN\ peaks with a significant rotation proxy it reaches \nmAucHarmDefined, but
over all peaks only \nmAucHarm, because the proxy is usually not detected: it is significant for \nmProxyActPct\% of the activity
impostors and \nmProxyPlanetPct\% of the planets alone. The number of
epochs alone has an AUC of \nmAucLogN: with both planet and activity amplitudes scaled to
the star's scatter, it carries no information about the class.

\begin{table}
\caption{Power of each ladder feature alone to separate planets from impostors
among \nmNDet\ significant simulated peaks.}
\label{tab:auc}
\centering
\footnotesize
\begin{tabular}{lc|lc}
\hline\hline
Feature & AUC & Feature & AUC \\
\hline
\inputrows{tab_auc_rows}
\hline
\end{tabular}
\tablefoot{AUC is reported as $\max(A, 1-A)$, so it does not depend on the sign of
the feature, and after the transforms of Sect.~\ref{sec:classifiers}. $^\dagger$Reported
but not used by the classifiers, because its power depends on the injection priors.}
\end{table}

\subsection{Classifiers}
\label{sec:classifiers}

We trained a logistic regression ($L_2$ penalty, $C = 0.5$, standardised features,
median imputation with missing-value indicators) and gradient-boosted trees
(scikit-learn's histogram implementation with 400 iterations, learning rate 0.04,
15 leaves, $L_2 = 1$ and at least 40 cases per leaf; \citealt{Pedregosa2011}) on the
ladder features of the significant peaks, excluding $\log N$ and $\log(T f_0)$. Heavy-tailed
features enter as $\log(1+x)$, the block probabilities as $\log_{10} p$, the alias
margin as ${\rm sign}(x)\log(1+|x|)$, and a missing \code{harm\_dist} as $\log(1+1000)$
with the flag \code{harm\_missing}. We used five-fold cross-validation in which all
cases of a star are in the same fold. Out of fold, the AUC is \nmAucLog\ for the
logistic regression and \nmAucTree\ for the trees (95\% intervals
\nmAucLogLo--\nmAucLogHi\ and \nmAucTreeLo--\nmAucTreeHi\ from a bootstrap over stars,
which holds the fitted models fixed), against \nmAucDetOnly\ for the detection features
alone (FAP, $\dchi$, $K/\sigma_K$) and \nmAucFapOnly\ for the FAP alone
(Fig.~\ref{fig:roc}). The combination, not any one test, carries the information. Adding
$\log N$ and $\log(T f_0)$ raises the trees' AUC to \nmAucTreePrior, through the period
distributions assumed for planets and rotation. Removing every activity-indicator and
rotation feature lowers it to \nmAucNoInd, and removing only the rotation-proxy features
to \nmAucNoProxy; the indicator power dominates the permutation importance
(Fig.~\ref{fig:importance}).

The Brier scores are \nmBrierLog\ and \nmBrierTree. Calibration was tested per bin of
predicted probability (Fig.~\ref{fig:reliability}): for the logistic regression the
largest deviation is $\nmCalMaxZLog\sigma$, with \nmCalNBadLog\ of \nmCalNBinsLog\ bins
beyond $2\sigma$, so it is well calibrated; the trees are overconfident at both ends, with
a largest deviation of $\nmCalMaxZTree\sigma$ and \nmCalNBadTree\ of \nmCalNBinsTree\ bins
beyond $2\sigma$. Either calibration holds only for the class mix of the simulations, so
on real candidates the output is a score, not a probability. At a
threshold of 0.5 the trees accept \nmTreeTPR\% of planets and \nmTreeFPR\% of
impostors: \nmTreeAccRotHarm\% at a rotation harmonic, \nmTreeAccRotLine\% inside a
rotation line, \nmTreeAccActOther\% of other activity, \nmTreeAccPlAlias\% of planet
aliases and \nmTreeAccNoise\% of noise peaks. At a false-positive rate of 5\% they
accept \nmTreeTPRatFiveFPR\% of planets, and at 1\%, \nmTreeTPRatOneFPR\%. If one
significant signal in ten were a planet, the precision at the 0.5 threshold would be
\nmTreePPVTen\%; at one in two, \nmTreePPVFifty\%.

\begin{figure}
\centering
\includegraphics[width=\columnwidth]{fig_roc.pdf}
\caption{Separating planets from impostors among significant simulated peaks, out of
fold with folds grouped by star; AUC in brackets. The square is the a-priori rule.}
\label{fig:roc}
\end{figure}

\begin{figure}
\centering
\includegraphics[width=\columnwidth]{fig_reliability.pdf}
\caption{Reliability of the out-of-fold probabilities: the fraction of cases that
are planets in bins of predicted probability, with binomial errors.}
\label{fig:reliability}
\end{figure}

\begin{figure}
\centering
\includegraphics[width=\columnwidth]{fig_importance.pdf}
\caption{Permutation importance of the features for the boosted trees, measured on
a quarter of the stars held out of training (symmetric-log scale).}
\label{fig:importance}
\end{figure}

\subsection{The a-priori rule}

The rule accepts \nmRuleTPR\% of planets and \nmRuleFPR\% of impostors; at the same
planet acceptance the trees accept \nmTreeFPRatRuleTPR\% of impostors (Fig.~\ref{fig:roc}). It accepts
\nmRuleAccRotHarm\% of impostors at a rotation harmonic, \nmRuleAccRotLine\% of those
inside a rotation line, \nmRuleAccActOther\% of other activity, \nmRuleAccPlAlias\% of
planet aliases and \nmRuleAccNoise\% of noise peaks. Its thresholds sit where they
reject few planets, and about half of the impostors in these simulations cross none of
them: the
indicators are weak or silent in many activity cases (Sect.~\ref{sec:sens}), a rotation
proxy is usually not detected, and a noise peak, once significant, looks stationary.

\subsection{Sensitivity to the simulated activity}
\label{sec:sens}

How much the indicators help depends on how strongly activity shows in them, which the
simulations set by assumption, and it is the largest dependence we find. For the \nmStNUnc\ activity impostors whose indicators
received nothing, the rule accepts \nmStRuleUnc\% and the trees \nmStTreeUnc\%; for
indicator amplitudes of 0.1--0.3 times the scatter (\nmStNWeak\ impostors), \nmStRuleWeak\%
and \nmStTreeWeak\%; for 1--3 times (\nmStNStrong), \nmStRuleStrong\% and \nmStTreeStrong\%.
The RV amplitude of the activity matters too: for amplitudes of 0.2--0.5 times the
star's scatter (\nmAmpLoN\ impostors) the rule accepts \nmAmpLoRule\% and the trees
\nmAmpLoTree\%, above once the scatter (\nmAmpHiN) \nmAmpHiRule\% and \nmAmpHiTree\%,
presumably because weak activity becomes significant mainly where it is coherent.
The rotation domain matters as well. For activity impostors with $P_{\rm rot} > 40$\,d
(\nmDomProtLongN), the rule accepts \nmDomProtLongRule\% and the trees
\nmDomProtLongTree\%; with $P_{\rm rot} < 10$\,d (\nmDomProtShortN), \nmDomProtShortRule\%
and \nmDomProtShortTree\%. For spot patterns that live longer than ten rotations
(\nmDomLamLongN), \nmDomLamLongRule\% and \nmDomLamLongTree\%; shorter than three
(\nmDomLamShortN), \nmDomLamShortRule\% and \nmDomLamShortTree\%. The \nmDomBothN\
impostors with both a long rotation period and long-lived spots, the regime of
GJ\,581\,d, are accepted at \nmDomBothRule\% by the rule and \nmDomBothTree\% by the trees,
against \nmRuleAccAct\% and \nmTreeAccAct\% for all activity impostors: over the ranges
simulated, the rotation period and the spot lifetime matter much less than the
indicators.
